% ============================================================================
%  SkillSync — Prompt Diary
%  From Vibe-Coded Prototype to a Production-Grade Smart Freelancing Platform
%
%  Compile with pdflatex (TeX Live / MiKTeX / Overleaf).
% ============================================================================
\documentclass[11pt,a4paper]{article}

\usepackage[T1]{fontenc}
\usepackage{lmodern}
\usepackage[margin=2.2cm]{geometry}
\usepackage{parskip}
\usepackage{xcolor}
\usepackage{booktabs}
\usepackage{tabularx}
\usepackage{enumitem}
\usepackage[most]{tcolorbox}
\usepackage{titlesec}
\usepackage[colorlinks=true,linkcolor=brandblue,urlcolor=brandblue]{hyperref}

% ---------------------------------------------------------------------------
% Brand colours
% ---------------------------------------------------------------------------
\definecolor{brandblue}{HTML}{1D4ED8}
\definecolor{brandink}{HTML}{0F172A}
\definecolor{brandmist}{HTML}{EEF2FF}
\definecolor{brandline}{HTML}{C7D2FE}
\definecolor{donegreen}{HTML}{047857}
\definecolor{phasegray}{HTML}{475569}

\titleformat{\section}{\Large\bfseries\color{brandink}}{\thesection}{0.8em}{}
\titleformat{\subsection}{\large\bfseries\color{phasegray}}{\thesubsection}{0.7em}{}

% ---------------------------------------------------------------------------
% Prompt box environment
%   \begin{promptbox}{<number>}{<title>}  ... \end{promptbox}
% ---------------------------------------------------------------------------
\newtcolorbox{promptbox}[2]{
  enhanced, breakable,
  colback=brandmist, colframe=brandblue,
  boxrule=0.8pt, arc=2mm,
  title={\textbf{Prompt #1 --- #2}},
  fonttitle=\bfseries\color{white},
  colbacktitle=brandblue,
  left=3mm, right=3mm, top=2.5mm, bottom=2.5mm
}

% "Definition of done" mini-heading
\newcommand{\dod}{\medskip\noindent\textcolor{donegreen}{\rule{\linewidth}{0.5pt}}\\[1mm]\textbf{\textcolor{donegreen}{Definition of done}}\\[-2mm]}
\newcommand{\objective}[1]{\noindent\textbf{Objective.} #1\medskip}
\newcommand{\filepath}[1]{\texttt{#1}}

% ============================================================================
\begin{document}

% ---------------------------------------------------------------------------
% Title page
% ---------------------------------------------------------------------------
\begin{titlepage}
  \centering
  \vspace*{2.5cm}
  {\Huge\bfseries\color{brandblue} SkillSync\par}
  \vspace{0.4cm}
  {\Large\color{brandink} Smart Freelancer Hiring \& Review System\par}
  \vspace{1.2cm}
  {\huge\bfseries The Prompt Diary\par}
  \vspace{0.5cm}
  {\large From Vibe-Coded Prototype to a Production-Grade\\Intelligent Freelancing Platform\par}
  \vspace{2cm}
  \begin{tcolorbox}[colback=brandmist,colframe=brandline,arc=2mm,width=0.8\textwidth]
    \centering
    \textbf{10 Phases \quad\textbullet\quad 27 Prompts \quad\textbullet\quad 1 Transformation}\\[2mm]
    Execute the prompts in order. Each prompt assumes every prompt before it\\has been completed and verified.
  \end{tcolorbox}
  \vfill
  {\color{phasegray} July 2026 \quad\textbullet\quad Version 1.0\par}
\end{titlepage}

\tableofcontents
\newpage

% ===========================================================================
\section{How to Use This Diary}
% ===========================================================================

This document is a \emph{development playbook}, not a code listing. Each prompt is a
self-contained instruction you paste into your AI coding session. Rules of engagement:

\begin{enumerate}[leftmargin=1.6em]
  \item \textbf{Order matters.} The sequence is deliberately dependency-ordered:
        identity $\rightarrow$ architecture $\rightarrow$ core loops $\rightarrow$
        intelligence $\rightarrow$ hardening. Skipping ahead will force rework.
  \item \textbf{One prompt per session.} Finish, verify against the ``Definition of
        done'' checklist, commit, and only then move on. Keep commits small and named
        after the prompt (e.g.\ \texttt{P07: profile \& settings}).
  \item \textbf{Verify before you trust.} After every prompt run
        \texttt{npm run typecheck} and click through the affected screens. The current
        codebase's biggest disease is \emph{mock success} --- forms that pretend to
        submit. Never accept a screen you have not exercised against the real API.
  \item \textbf{Never regress the loop.} From Phase~3 onward there is always a working
        end-to-end hiring loop. Every later prompt must preserve it.
  \item \textbf{Adapt freely.} If your session discovers that a referenced file has
        drifted, tell the agent to re-inspect and adapt --- the \emph{intent} of each
        prompt is the contract, not the exact file list.
\end{enumerate}

% ===========================================================================
\section{State of the Project (Condensed Audit)}
% ===========================================================================

A full audit of the repository (frontend, backend, models, docs) produced the
following verdict. The backend is a respectable skeleton; the frontend is a
marketing shell with an almost entirely simulated application behind it.

\subsection{What already works}
\begin{itemize}[leftmargin=1.6em]
  \item $\sim$34 real API routes: auth (credentials + Google, password reset),
        projects, skills, proposals (with a proper MongoDB transaction on accept),
        contracts, Stripe PaymentIntent + verified webhook, reviews (race-safe
        uniqueness), notifications, admin user management, search with text index.
  \item Sound security hygiene in places: bcrypt cost 12, hashed reset tokens,
        \texttt{select:\,false} on secrets, enumeration-safe responses,
        denormalised rating aggregates.
  \item A polished landing-page aesthetic: animations, dark mode, glassmorphism.
\end{itemize}

\subsection{Critical defects}
\begin{itemize}[leftmargin=1.6em]
  \item \textbf{Identity crisis.} Skill-\emph{learning} language (``Learn Skills'',
        ``Teach \& Mentor'', ``Learner''/``Expert'') collides with freelancing
        mechanics (proposals, contracts). Roles are named \texttt{member}/\texttt{provider}
        instead of client/freelancer. Two clashing colour systems (blue/emerald vs.\
        indigo/zinc). Two contradictory category taxonomies.
  \item \textbf{No marketing/app separation.} One global layout wraps everything;
        login redirects to the landing page; logged-in users keep marketing navigation.
  \item \textbf{Simulated frontend.} Post-Project, Share-Skill and Hire dialogs fake
        success with \texttt{setTimeout}; the dashboard shows hardcoded stats and a
        fake activity feed; \texttt{/hire-talent} renders 48 randomly generated
        freelancers; every profile is a hardcoded ``Alex Morgan''; the Stripe checkout
        component is imported by no page.
  \item \textbf{Ghost routes.} \texttt{/register}, \texttt{/explore},
        \texttt{/profile}, \texttt{/settings}, \texttt{/projects/[id]},
        \texttt{/dashboard/contracts/[id]} are linked or documented but do not exist.
        Freelancers cannot browse or bid on projects at all.
  \item \textbf{Broken economics.} No Stripe Connect, no payout path, no escrow
        despite the Terms page promising it; a freelancer can self-complete an unpaid
        contract; cancelled contracts can still be funded.
  \item \textbf{Enforcement gaps.} Banned users can log in and use every API; JWT
        roles go stale after admin changes; $\sim$9 routes return 500 on malformed
        ObjectIds; no rate limit on login or forgot-password; unescaped user-supplied
        regex in four routes.
  \item \textbf{Dead code.} \texttt{Booking} model, all of
        \filepath{lib/aggregations.ts}, unused helpers in \filepath{lib/api-utils.ts},
        seed and index scripts that write or index phantom fields.
  \item \textbf{Zero AI.} ``AI matching'' exists only as marketing copy.
\end{itemize}

\subsection{Gap matrix vs.\ industry platforms}
\begin{table}[h]
\small
\begin{tabularx}{\textwidth}{@{}lXl@{}}
\toprule
\textbf{Capability} & \textbf{Current state} & \textbf{Phase} \\
\midrule
Landing / app split          & None --- one shared shell            & 1 \\
Role-based onboarding        & Single radio button, wrong labels    & 2 \\
Client dashboard             & Mock data, dead tabs                 & 3 \\
Project browse \& bidding    & Missing entirely                     & 4 \\
Freelancer dashboard         & Mock data, dead tabs                 & 4 \\
Contract workspace           & Missing (API exists)                 & 5 \\
Escrow / payouts / earnings  & Missing (Terms page lies)            & 5 \\
Messaging                    & Missing entirely                     & 6 \\
Notifications UI             & Missing (API exists)                 & 6 \\
Reviews UI                   & Missing (API exists)                 & 6 \\
AI features                  & Missing entirely                     & 7 \\
Admin panel                  & Missing (3 API routes exist)         & 8 \\
Security hardening           & Partial                              & 9 \\
Tests                        & 0\,\% coverage                       & 9 \\
\bottomrule
\end{tabularx}
\end{table}

% ===========================================================================
\section{Target Architecture}
% ===========================================================================

\subsection{Naming and roles}
The platform's vocabulary is unified once, in Phase 0, and never violated again:

\begin{itemize}[leftmargin=1.6em]
  \item \texttt{member} $\rightarrow$ \textbf{client} (hires talent);
        \texttt{provider} $\rightarrow$ \textbf{freelancer} (provides services);
        \texttt{admin} stays.
  \item ``Skill'' listings become \textbf{Gigs} (Fiverr-style service offerings).
  \item All learning/teaching/mentoring language is eliminated.
\end{itemize}

\subsection{Route topology}
\begin{itemize}[leftmargin=1.6em]
  \item \texttt{(marketing)} --- landing, about, contact, terms, privacy. Public.
        Marketing navigation and footer. Authenticated visitors hitting \texttt{/}
        are redirected into the app.
  \item \texttt{(auth)} --- login, register, onboarding, password reset. Minimal
        chrome, no marketing navigation.
  \item \texttt{(app)} --- everything behind login. Its own shell: role-aware
        sidebar, topbar with notifications and messages, no marketing links.
        Client area under \texttt{/client/\ldots}, freelancer area under
        \texttt{/freelancer/\ldots}, shared surfaces (contracts, messages, settings)
        role-agnostic.
  \item \texttt{(admin)} --- \texttt{/admin/\ldots}, gated by role at middleware level.
\end{itemize}

\subsection{The Smart layer}
All AI lives server-side behind \filepath{lib/ai/} using the Gemini API
(\texttt{GEMINI\_API\_KEY} in env; never exposed to the client). Six pillars:

\begin{enumerate}[leftmargin=1.6em]
  \item \textbf{Smart Matching Engine} --- Gemini embeddings over project briefs and
        freelancer profiles; ranked matches with human-readable reasons.
  \item \textbf{Hiring Copilot} --- proposal ranking, summarisation and red-flag
        detection for clients.
  \item \textbf{AI Writing Studio} --- job-post composer, proposal assistant,
        profile optimiser.
  \item \textbf{Smart Pricing} --- budget estimation for clients, rate guidance for
        freelancers.
  \item \textbf{Review Intelligence} --- sentiment digests on profiles, anomalous
        review flagging.
  \item \textbf{SyncMate} --- a role-aware platform assistant widget.
\end{enumerate}

Every AI feature degrades gracefully: if the Gemini call fails or the key is absent,
the surrounding feature still works without its smart enhancement.

\newpage
% ===========================================================================
\section{Phase 0 --- Foundation Reset}
% ===========================================================================

\textit{Purpose: kill the identity crisis and the dead weight before building
anything new. Everything later assumes this vocabulary and hygiene.}

% ---------------------------------------------------------------------------
\objective{Rename the roles across the entire stack, unify all product language
around freelancing, and collapse the two category taxonomies into one.}

\begin{promptbox}{01}{Product Identity \& Role Unification}
You are working on SkillSync, a Next.js 15 App Router + Mongoose + NextAuth v5
freelancing platform. Before adding features, unify its identity. Do the following
across the entire codebase:

\begin{enumerate}[leftmargin=1.6em]
  \item Rename the user roles: \texttt{member} $\rightarrow$ \texttt{client} and
        \texttt{provider} $\rightarrow$ \texttt{freelancer} (keep \texttt{admin}).
        Update the Mongoose \texttt{User} schema enum, every Zod schema, every TS type,
        the NextAuth JWT/session callbacks in \filepath{lib/auth.config.ts} and
        \filepath{lib/auth.ts}, every API route check, and every UI conditional.
        Write a one-off migration script \filepath{scripts/migrate-roles.ts} that
        updates existing documents in MongoDB.
  \item Eliminate ALL skill-learning language. Replace ``Learn Skills / Teach \&
        Mentor'' on the signup page with ``I'm a Client --- I want to hire talent''
        and ``I'm a Freelancer --- I want to offer services''. Replace
        ``Learner''/``Expert'' badges in the dashboard. Rewrite the hero headline,
        metadata tagline (``Share skills. Hire talent.''), and any copy mentioning
        learning, teaching, mentoring, or lessons so everything reads as a modern
        freelancer hiring platform.
  \item Merge the two conflicting category taxonomies (\filepath{data/categories.ts}
        vs.\ the lists inside hire-talent and post-project) into ONE canonical list in
        \filepath{data/categories.ts} with freelancing categories only (Web
        Development, Mobile, UI/UX Design, Graphic Design, Content Writing, Digital
        Marketing, Video \& Animation, Data \& AI, DevOps \& Cloud, Business \&
        Consulting). Import it everywhere; delete the local copies.
  \item Pick ONE design token system --- the existing brand blue/emerald oklch tokens
        in \filepath{app/globals.css} --- and replace all hardcoded indigo/purple/zinc
        utility classes in the dashboard, auth, hire-talent, freelancer profile, and
        form pages with the semantic tokens.
  \item Do not add new features in this prompt. Run \texttt{npm run typecheck} and fix
        every error you introduced.
\end{enumerate}

\dod
\begin{itemize}[leftmargin=1.6em]
  \item \texttt{grep} for ``member''/``provider'' role strings, ``Learner'',
        ``Expert'', ``Teach'', ``mentor'' returns no product-facing hits.
  \item Signup shows Client vs Freelancer cards; typecheck passes; app boots.
  \item One category list; one colour system.
\end{itemize}
\end{promptbox}

% ---------------------------------------------------------------------------
\objective{Remove every piece of dead or lying code so later prompts build on
truth, and repair the data layer scripts.}

\begin{promptbox}{02}{Dead Code Purge \& Data-Layer Repair}
Clean the SkillSync codebase of dead and inconsistent code. Do not add features.

\begin{enumerate}[leftmargin=1.6em]
  \item Delete the unused \filepath{models/Booking.ts}, its \texttt{Booking} types in
        \filepath{types/index.ts}, its Zod schema in \filepath{types/schemas.ts}, and
        the \texttt{/bookings} entry in \filepath{middleware.ts}.
  \item Delete \filepath{lib/aggregations.ts} entirely (all four exports are unused
        and reference a non-existent \texttt{isVerified} field). Remove the unused
        \texttt{withErrorHandler} and \texttt{paginatedResponse} from
        \filepath{lib/api-utils.ts}, and the dead first \texttt{pipeline} variable in
        \filepath{app/api/proposals/received/route.ts}.
  \item Add the missing \texttt{hourlyRate} (number, min 0) and \texttt{title}-style
        professional fields the profile API already accepts to the \texttt{User}
        schema so \texttt{PUT /api/users/me} stops silently dropping them.
  \item Fix \filepath{scripts/create-indexes.ts}: remove indexes on non-existent
        fields (\texttt{isVerified}, \texttt{availability}, \texttt{isAvailable},
        \texttt{rate}) and index the fields that actually exist.
  \item Rewrite \filepath{scripts/seed.ts} so it only writes fields that exist in the
        schemas, uses the new \texttt{client}/\texttt{freelancer} roles and the
        canonical categories, calls the rating recalculation in
        \filepath{lib/reviews.ts} after seeding reviews, and produces a coherent demo
        dataset: $\sim$12 freelancers with rich profiles, $\sim$6 clients,
        $\sim$15 projects in various states, proposals, a few contracts and reviews.
  \item Remove the no-op Pages-Router \texttt{config} export from
        \filepath{app/api/payments/webhook/route.ts}.
  \item Update \filepath{docs/} files only where they now state falsehoods (roles,
        deleted files). Run \texttt{npm run typecheck}; re-run the seed against the dev
        database to prove it works.
\end{enumerate}

\dod
\begin{itemize}[leftmargin=1.6em]
  \item No references to Booking, aggregations.ts, or phantom fields remain.
  \item \texttt{npm run seed} completes and seeded users show correct average ratings.
\end{itemize}
\end{promptbox}

% ===========================================================================
\section{Phase 1 --- Architecture Split \& The Landing Page}
% ===========================================================================

\textit{Purpose: create the marketing/app separation the product demands ---
a public showcase site and, behind login, a real application shell.}

% ---------------------------------------------------------------------------
\objective{Restructure the route tree into route groups with independent layouts,
and fix every broken link in the codebase.}

\begin{promptbox}{03}{Route Groups, Layouts \& Link Integrity}
Restructure SkillSync's App Router tree into three route groups with separate
layouts. Do not redesign pages yet --- move and wire only.

\begin{enumerate}[leftmargin=1.6em]
  \item Create \texttt{app/(marketing)/} containing the landing page \texttt{/},
        about, services, contact, terms, privacy. Its layout renders the existing
        marketing \texttt{Navigation} and \texttt{Footer}.
  \item Create \texttt{app/(auth)/} containing login, register (rename signup
        $\rightarrow$ register so the nav's ``Get Started'' link finally works; keep a
        redirect from \texttt{/signup}), forgot-password, reset-password. Its layout
        is the minimal split-screen \texttt{AuthLayout} with NO marketing nav or
        footer.
  \item Create \texttt{app/(app)/} for everything behind login: move
        \texttt{/dashboard}, \texttt{/post-project}, \texttt{/hire-talent},
        \texttt{/share-skill} here. Give it a placeholder application layout (a plain
        sidebar + topbar scaffold; Prompt 05 will design it properly). No marketing
        nav or footer inside.
  \item Middleware: protect everything in the app group; when an authenticated user
        visits \texttt{/}, \texttt{/login}, or \texttt{/register}, redirect them to
        \texttt{/dashboard}. Login must redirect to \texttt{/dashboard} (honouring
        \texttt{callbackUrl}), never to \texttt{/}. Keep \texttt{/hire-talent} inside
        the authenticated app group for now; public talent discovery can be added as
        a marketing surface later if desired.
  \item Fix every broken link: navigation ``Get Started'' $\rightarrow$
        \texttt{/register}; footer links to \texttt{/explore} and
        \texttt{/projects/new} $\rightarrow$ real routes; user-menu
        \texttt{/profile} and \texttt{/settings} $\rightarrow$ keep the links but
        create minimal placeholder pages so nothing 404s; \texttt{ProtectedRoute}'s
        \texttt{/unauthorized} $\rightarrow$ create a simple unauthorized page;
        remove links to routes that no longer exist.
  \item Add \filepath{app/not-found.tsx}, \filepath{app/error.tsx} and
        \filepath{app/global-error.tsx} with branded, friendly fallbacks.
  \item Verify by clicking every link in nav, footer, user menu, and dashboard quick
        actions --- zero 404s. Typecheck passes.
\end{enumerate}

\dod
\begin{itemize}[leftmargin=1.6em]
  \item Three route groups with independent layouts; login lands on
        \texttt{/dashboard}; logged-in users never see marketing chrome.
  \item Zero dead links anywhere; custom 404/error pages exist.
\end{itemize}
\end{promptbox}

% ---------------------------------------------------------------------------
\objective{Rebuild the landing page as a credible, freelancing-focused product
showcase --- the professional first impression.}

\begin{promptbox}{04}{The Landing Page}
Redesign the SkillSync landing page in \texttt{app/(marketing)/} as a
production-quality marketing site for a SMART freelancing platform. Keep the
existing design tokens (brand blue/emerald oklch system) and animation utilities.
Sections, in order:

\begin{enumerate}[leftmargin=1.6em]
  \item \textbf{Hero}: headline positioning SkillSync as AI-powered freelancer
        hiring (e.g.\ ``Hire smarter. Work smarter.''), subline explaining the
        platform in one sentence, dual CTAs (``Hire a Freelancer'' / ``Become a
        Freelancer'' --- both to \texttt{/register} with a \texttt{?role=} hint), and
        a visual that suggests intelligent matching. Remove the fabricated
        ``50K+ users'' stats or clearly label the section as illustrative.
  \item \textbf{How it works}: two tracks side by side --- For Clients (post
        $\rightarrow$ get AI-matched proposals $\rightarrow$ hire \& pay securely)
        and For Freelancers (create profile $\rightarrow$ get matched to projects
        $\rightarrow$ deliver \& get paid).
  \item \textbf{Smart features grid}: showcase the six AI pillars --- Smart Matching,
        Hiring Copilot, AI Writing Studio, Smart Pricing, Review Intelligence,
        SyncMate assistant --- as the differentiator. Only describe what will exist.
  \item \textbf{Category strip} using the canonical categories.
  \item \textbf{Trust section}: secure payments, verified reviews, dispute support.
  \item \textbf{About section} (brief, on-page) + final CTA + existing footer,
        with all links valid.
\end{enumerate}
The page must be a Server Component except leaf interactive islands, fully
responsive, dark-mode correct, and free of any learning/teaching language.

\dod
\begin{itemize}[leftmargin=1.6em]
  \item Landing reads unmistakably as an AI-powered freelancing platform.
  \item Lighthouse-reasonable: no layout shift, images optimised, responsive at
        375px/768px/1280px.
\end{itemize}
\end{promptbox}

% ---------------------------------------------------------------------------
\objective{Build the real application shell every logged-in user lives in.}

\begin{promptbox}{05}{The Application Shell}
Design the authenticated application shell for \texttt{app/(app)/layout.tsx},
replacing the placeholder from Prompt 03. This is the workspace chrome for both
roles --- model it on the information density of Upwork's app, not a marketing site.

\begin{enumerate}[leftmargin=1.6em]
  \item \textbf{Sidebar} (collapsible on mobile): logo; role-aware primary nav.
        Client items: Dashboard, My Projects, Find Freelancers, Proposals Received,
        Contracts, Messages, Payments. Freelancer items: Dashboard, Find Work, My
        Gigs, My Proposals, Contracts, Messages, Earnings. Shared bottom items:
        Settings, Help. Highlight the active route.
  \item \textbf{Topbar}: global search placeholder (wired later), notifications bell
        with unread badge fed by \texttt{GET /api/notifications} (poll every 60s),
        theme toggle, avatar menu (Profile, Settings, Log out).
  \item Notifications bell opens a dropdown listing recent notifications with
        read/unread state, wired to the existing mark-read APIs, plus a ``View all''
        link to a \texttt{/notifications} page (simple list with pagination and
        ``mark all read'').
  \item Routes that do not exist yet should render tidy ``coming soon'' placeholders
        within the shell --- never 404 --- so navigation is complete from day one.
  \item Ensure the shell renders user data from the session server-side (no flash of
        unauthenticated content), works at mobile widths with a drawer, and respects
        the design tokens.
\end{enumerate}

\dod
\begin{itemize}[leftmargin=1.6em]
  \item Both roles see coherent, distinct navigation; bell shows real unread counts;
        every sidebar item routes somewhere sensible.
\end{itemize}
\end{promptbox}

% ===========================================================================
\section{Phase 2 --- Auth \& Onboarding}
% ===========================================================================

\textit{Purpose: turn account creation into a proper role-driven onboarding
journey that lands each user in the right dashboard with a usable profile.}

% ---------------------------------------------------------------------------
\objective{A multi-step onboarding that makes role selection meaningful and
eliminates the ``register, then log in again manually'' friction.}

\begin{promptbox}{06}{Registration \& Role-Based Onboarding Wizard}
Rebuild registration as a multi-step onboarding flow in \texttt{app/(auth)/register}.

\begin{enumerate}[leftmargin=1.6em]
  \item \textbf{Step 1 --- Role selection}: full-screen choice between two rich cards:
        ``I'm a Client --- I want to hire'' and ``I'm a Freelancer --- I want to
        work''. Honour a \texttt{?role=} query param from the landing CTAs as
        preselection.
  \item \textbf{Step 2 --- Account}: name, email, password with live strength
        criteria, terms. On submit call \texttt{POST /api/auth/register}, then
        automatically sign the user in with the credentials (no manual re-login) and
        continue to step 3.
  \item \textbf{Step 3 --- Profile setup} (role-specific, saved via
        \texttt{PUT /api/users/me}): freelancers provide headline, hourly rate,
        skills (tag input), bio, and optional avatar upload; clients provide company/
        display name, what they typically hire for (categories), and optional avatar.
        Include a visible ``Skip for now'' escape hatch.
  \item \textbf{Step 4 --- Done}: short success screen, then redirect to
        \texttt{/dashboard}.
  \item \textbf{Google OAuth}: after first-time OAuth sign-in, users have no chosen
        role --- detect this and route them into the wizard at Step 1 (role +
        profile only). Persist their choice via a new safe API mechanism; the role
        must be settable exactly once by the user, then only by admins.
  \item Show a progress indicator across steps; all state survives refresh (persist
        step data client-side until committed).
\end{enumerate}

\dod
\begin{itemize}[leftmargin=1.6em]
  \item New credential users go card $\rightarrow$ form $\rightarrow$ profile
        $\rightarrow$ dashboard without ever re-entering credentials.
  \item OAuth first-timers are forced through role selection; role is immutable
        afterwards for non-admins.
\end{itemize}
\end{promptbox}

% ---------------------------------------------------------------------------
\objective{Give the profile and settings links real destinations.}

\begin{promptbox}{07}{Profile Management \& Settings}
Build the profile and settings surfaces inside the app shell.

\begin{enumerate}[leftmargin=1.6em]
  \item \texttt{/settings/profile}: full profile editor wired to
        \texttt{GET/PUT /api/users/me} --- avatar upload via the existing
        \texttt{/api/upload} route with client-side preview, name, headline, bio,
        location, hourly rate (freelancers), skills tag editor. Server-validated,
        optimistic UI with toasts.
  \item \texttt{/settings/account}: change password (add the missing API route:
        verify current password, validate new one with the shared Zod rules),
        connected sign-in methods display, and a danger zone placeholder for account
        deletion (UI only, disabled, labelled ``contact support'').
  \item \texttt{/settings/notifications}: preference toggles persisted on the User
        document (email on proposal, contract, payment, review events). The email
        helpers in \filepath{lib/email.ts} must respect these flags.
  \item Public freelancer profile page \texttt{/freelancers/[id]}: REPLACE the
        hardcoded ``Alex Morgan'' mock with real data from
        \texttt{GET /api/users/[id]} and \texttt{GET /api/reviews/[id]} --- header
        (avatar, name, headline, rate, location, rating), about, skills, gigs
        (from the skills API, by provider), and reviews with rating breakdown. Handle
        not-found IDs gracefully. Increment \texttt{profileViews} server-side on
        view (add the missing write path, excluding self-views).
\end{enumerate}

\dod
\begin{itemize}[leftmargin=1.6em]
  \item Editing a profile visibly changes the public profile page.
  \item No mock person remains anywhere; profile views actually increment.
\end{itemize}
\end{promptbox}

% ===========================================================================
\section{Phase 3 --- The Client Experience}
% ===========================================================================

\textit{Purpose: make the client side real --- posting projects that persist,
and a dashboard that tells the truth.}

% ---------------------------------------------------------------------------
\objective{Wire project posting to the real API and give clients management of
their projects.}

\begin{promptbox}{08}{Real Project Posting \& Project Management}
Make the client project workflow real.

\begin{enumerate}[leftmargin=1.6em]
  \item \filepath{app/(app)/post-project/PostProjectForm.tsx} currently fakes
        submission with \texttt{setTimeout}. Wire it to
        \texttt{POST /api/projects} for real: map every wizard field to the API
        schema, upload attachments through \texttt{/api/upload} first and submit the
        returned URLs, handle validation errors field-by-field, and on success
        redirect to the new project's detail page. Add a cross-field check that
        \texttt{budgetMin} $\le$ \texttt{budgetMax} on both client and API.
  \item Build \texttt{/client/projects}: the client's project list (wired to
        \texttt{GET /api/projects} filtered to own), with status chips (open,
        in progress, completed, cancelled), proposal counts, and actions: view, edit
        (open projects only), close.
  \item Build the project detail page \texttt{/projects/[id]} (shared route, both
        roles can view): full brief, budget, skills, attachments, client card, and a
        contextual action area --- owner sees proposals summary and manage buttons;
        freelancers see a ``Submit Proposal'' CTA (wired in Prompt 11); visitors of
        the wrong state see status-appropriate messaging.
  \item Restrict \texttt{POST /api/projects} to the \texttt{client} role, and
        constrain the owner's status edits in \texttt{PUT /api/projects/[id]} to
        legal transitions only (open $\rightarrow$ cancelled; everything else is
        driven by the system).
\end{enumerate}

\dod
\begin{itemize}[leftmargin=1.6em]
  \item Posting a project stores a real document and appears in the client's list and
        on its detail page after refresh.
  \item Freelancer accounts cannot post projects.
\end{itemize}
\end{promptbox}

% ---------------------------------------------------------------------------
\objective{A client dashboard fed entirely by real data.}

\begin{promptbox}{09}{Client Dashboard with Real Data}
Replace the mock client dashboard. Delete the hardcoded stats, sparkline array and
fake activity feed in \filepath{DashboardClient.tsx} for the client role and build a
real \texttt{/dashboard} (client variant):

\begin{enumerate}[leftmargin=1.6em]
  \item KPI cards wired to \texttt{GET /api/dashboard/stats}: active projects,
        proposals awaiting review, active contracts, total spent. Extend that API if a
        number it needs is missing --- never invent numbers client-side.
  \item ``Needs your attention'' panel: newest pending proposals on your projects
        (from \texttt{GET /api/proposals/received}) with quick links.
  \item Recent activity: derive a real feed from the user's notifications.
  \item Active contracts strip with payment status chips.
  \item Quick actions: Post a Project, Find Freelancers, View Contracts --- all to
        real routes.
  \item Empty states for brand-new clients (friendly illustrations + primary CTA to
        post the first project). Loading skeletons for every async block.
\end{enumerate}
The dashboard page must be a Server Component fetching initial data server-side,
with client islands only where interactivity demands it.

\dod
\begin{itemize}[leftmargin=1.6em]
  \item A fresh client sees zeros and empty states; after posting a project and
        receiving a proposal, every number changes accordingly.
  \item No hardcoded statistics remain for the client role.
\end{itemize}
\end{promptbox}

% ===========================================================================
\section{Phase 4 --- The Freelancer Experience}
% ===========================================================================

\textit{Purpose: freelancers currently cannot even see projects. Build discovery,
bidding, gigs, and a truthful freelancer dashboard.}

% ---------------------------------------------------------------------------
\objective{The missing core: freelancers browsing and finding work.}

\begin{promptbox}{10}{Find Work --- Project Discovery}
Build \texttt{/freelancer/find-work}, the freelancer's project discovery surface
(the platform's most glaring absence).

\begin{enumerate}[leftmargin=1.6em]
  \item Card list of open projects from \texttt{GET /api/projects?status=open}:
        title, truncated description, budget display (fixed range vs.\ hourly),
        required skills, experience level, posted-ago, proposal count, client rating.
  \item Filter sidebar: keyword (server-side search), category, budget type, budget
        range, experience level; sort by newest / budget. Filters live in the URL
        query string so results are shareable and refresh-safe.
  \item Proper pagination wired to the API's existing pagination.
  \item Each card links to the shared \texttt{/projects/[id]} detail page from
        Prompt 08, where the ``Submit Proposal'' CTA lives.
  \item Escape all user-supplied search input before it reaches MongoDB regex/text
        filters (fix the unescaped \texttt{\$regex} usage in the projects, skills and
        freelancers routes while you are here).
  \item Also connect \texttt{/hire-talent} (rename to \texttt{/client/find-freelancers})
        to the REAL \texttt{GET /api/freelancers} API --- delete the 48 randomly
        generated mock freelancers --- with its filters mapped to real query params,
        and make every card link to the real \texttt{/freelancers/[id]} profile.
\end{enumerate}

\dod
\begin{itemize}[leftmargin=1.6em]
  \item Seeded projects appear with working filters; zero generated mock data
        remains on either discovery surface.
\end{itemize}
\end{promptbox}

% ---------------------------------------------------------------------------
\objective{Real bidding: submit, track, withdraw.}

\begin{promptbox}{11}{Proposals --- Submit \& Manage}
Implement the full proposal experience for freelancers.

\begin{enumerate}[leftmargin=1.6em]
  \item On \texttt{/projects/[id]}, build the proposal composer (modal or inline
        section): cover-letter message, proposed rate (prefilled from profile
        hourly rate for hourly projects), estimated timeline. Submit to the real
        \texttt{POST /api/proposals}; surface duplicate-bid and own-project errors
        cleanly; on success show the submitted state on the page.
  \item Restrict \texttt{POST /api/proposals} to the \texttt{freelancer} role.
  \item Build \texttt{/freelancer/proposals}: tabbed list (pending / accepted /
        rejected / withdrawn) from \texttt{GET /api/proposals/sent}, each row linking
        to its project, showing status changes over time.
  \item Add the missing withdraw capability: a new authorised state transition
        allowing the proposal's own freelancer to set \texttt{pending}
        $\rightarrow$ \texttt{withdrawn} (extend \texttt{PUT /api/proposals/[id]}
        with an ownership branch), with a confirm dialog in the UI.
  \item Emit the currently-missing in-app notifications: notify the project owner on
        new proposal (\texttt{proposal\_received}) and the freelancer on
        accept/reject (\texttt{proposal\_update}), alongside the existing emails.
\end{enumerate}

\dod
\begin{itemize}[leftmargin=1.6em]
  \item A freelancer can bid on a seeded project, see it under Proposals, withdraw
        it, and the client receives an in-app notification for the bid.
  \item Clients cannot submit proposals.
\end{itemize}
\end{promptbox}

% ---------------------------------------------------------------------------
\objective{Turn ``skills'' into proper Fiverr-style gigs with public pages.}

\begin{promptbox}{12}{Gigs --- Service Offerings}
Convert the ``Share Skill'' feature into a proper Gigs system.

\begin{enumerate}[leftmargin=1.6em]
  \item Rename the surface: \texttt{/share-skill} $\rightarrow$
        \texttt{/freelancer/gigs/new}; all ``skill listing'' copy becomes ``gig''.
        Keep the \texttt{Skill} model name internally if renaming is disruptive, but
        every user-facing string says Gig.
  \item Wire the creation form (currently a \texttt{setTimeout} mock) to the real
        \texttt{POST /api/skills}, including category (canonical list), delivery
        expectations and tags.
  \item Build \texttt{/freelancer/gigs}: the freelancer's own gig list from
        \texttt{GET /api/skills/mine} with publish/unpublish toggle (wired to the
        \texttt{isPublished} field via \texttt{PUT /api/skills/[id]}), edit and
        delete (existing APIs).
  \item Build the public gig detail page \texttt{/gigs/[slug]}: gig content, the
        freelancer card with rating, similar gigs strip, and a ``Contact'' /
        ``Invite to project'' CTA for clients (contact wiring arrives with
        messaging; render it disabled with a tooltip until then).
  \item Show a freelancer's published gigs on their public profile page.
\end{enumerate}

\dod
\begin{itemize}[leftmargin=1.6em]
  \item Creating a gig persists it; it appears in My Gigs, on the public profile,
        and at its own URL. Unpublishing hides it publicly.
\end{itemize}
\end{promptbox}

% ---------------------------------------------------------------------------
\objective{A freelancer dashboard that reports reality.}

\begin{promptbox}{13}{Freelancer Dashboard with Real Data}
Build the freelancer variant of \texttt{/dashboard} with real data, mirroring the
quality bar of Prompt 09:

\begin{enumerate}[leftmargin=1.6em]
  \item KPI cards from \texttt{GET /api/dashboard/stats}: active contracts, pending
        proposals, total earned (completed contracts), average rating. Extend the
        stats API where needed.
  \item Proposal pipeline widget: counts by status with a link to
        \texttt{/freelancer/proposals}.
  \item ``Recommended for you'' placeholder panel querying newest open projects in
        the freelancer's categories (this becomes AI-powered in Phase 7 --- structure
        it so the data source can be swapped).
  \item Earnings snapshot: simple monthly bar chart of completed contract values
        (real aggregation, added to the stats API).
  \item Recent reviews received; recent activity from notifications; empty states
        and skeletons throughout.
  \item Delete every remaining hardcoded stat, sparkline, and fake activity item
        from the old dashboard for both roles; remove the dead sidebar-tab state
        machine.
\end{enumerate}

\dod
\begin{itemize}[leftmargin=1.6em]
  \item Zero mock data remains in any dashboard; numbers move when seeded data
        changes.
\end{itemize}
\end{promptbox}

% ===========================================================================
\section{Phase 5 --- Hiring, Contracts \& Money}
% ===========================================================================

\textit{Purpose: close the loop --- proposals get reviewed, contracts get a real
workspace, and money flows with escrow semantics the Terms page already promises.}

% ---------------------------------------------------------------------------
\objective{The client's hiring decision surface.}

\begin{promptbox}{14}{Proposal Review \& Hiring}
Build the client's proposal-review experience.

\begin{enumerate}[leftmargin=1.6em]
  \item \texttt{/client/proposals}: all incoming proposals across the client's
        projects from \texttt{GET /api/proposals/received}, grouped by project,
        filterable by status.
  \item On the owner's view of \texttt{/projects/[id]}: a proposals tab listing each
        bid --- freelancer card (avatar, rating, rate, headline linking to profile),
        cover letter, proposed rate and timeline --- with Accept and Reject actions
        wired to the existing transactional \texttt{PUT /api/proposals/[id]}.
  \item Accepting shows a confirmation dialog explaining the consequences (other
        pending proposals auto-rejected, contract created, project moves to in
        progress) and on success routes the client to the new contract page.
  \item Rejecting asks for an optional reason (stored on the proposal; extend the
        model with a nullable \texttt{clientNote}) and notifies the freelancer.
  \item Handle the race where the project is no longer open gracefully.
\end{enumerate}

\dod
\begin{itemize}[leftmargin=1.6em]
  \item Accept creates the contract and auto-rejects siblings (verify in DB);
        both parties receive notifications.
\end{itemize}
\end{promptbox}

% ---------------------------------------------------------------------------
\objective{A shared contract workspace where the funded work actually lives.}

\begin{promptbox}{15}{Contract Workspace}
Build \texttt{/contracts} (list, both roles) and \texttt{/contracts/[id]} (workspace),
finally putting the orphaned Stripe \texttt{CheckoutForm} to use.

\begin{enumerate}[leftmargin=1.6em]
  \item List page: the caller's contracts from \texttt{GET /api/contracts} with
        role-aware framing (client sees freelancer, freelancer sees client), status
        and payment chips.
  \item Workspace page from \texttt{GET /api/contracts/[id]}: header (project title,
        parties, agreed rate, timeline, started date), a status timeline
        (created $\rightarrow$ funded $\rightarrow$ delivered $\rightarrow$ completed),
        and a role-aware action panel.
  \item \textbf{Funding}: when payment is pending, the client sees ``Fund this
        contract'' opening the existing \texttt{CheckoutForm} wired to
        \texttt{POST /api/payments/create-intent}; handle success/failure states and
        reflect webhook-driven \texttt{paid} status (poll or refresh affordance).
  \item \textbf{Delivery flow}: add a \texttt{delivered} step --- the freelancer marks
        work delivered (new state on the contract); only then can the CLIENT mark the
        contract \texttt{completed}. Remove the current flaw where either party
        (including the freelancer) can unilaterally complete. Completion requires
        \texttt{paymentStatus === "paid"} --- enforce in
        \texttt{PUT /api/contracts/[id]}.
  \item \textbf{Guards}: \texttt{create-intent} must refuse non-active contracts;
        add an idempotency key; validate ObjectIds.
  \item Either party may raise a dispute (status \texttt{disputed}) with a required
        reason stored on the contract; disputes freeze other transitions until admin
        resolution (Phase 8).
  \item Full notification + email coverage for fund, deliver, complete, dispute.
\end{enumerate}

\dod
\begin{itemize}[leftmargin=1.6em]
  \item End-to-end with Stripe test cards: accept $\rightarrow$ fund $\rightarrow$
        deliver $\rightarrow$ complete works; freelancer cannot self-complete;
        unpaid contracts cannot complete.
\end{itemize}
\end{promptbox}

% ---------------------------------------------------------------------------
\objective{Honest money semantics: an earnings ledger and escrow-style release.}

\begin{promptbox}{16}{Earnings, Escrow Semantics \& Payment Hardening}
Give SkillSync truthful money mechanics on top of the single-payment flow.

\begin{enumerate}[leftmargin=1.6em]
  \item Create a \texttt{Transaction} model: contractId, clientId, freelancerId,
        amount, type (\texttt{funding}, \texttt{release}, \texttt{refund}), status,
        stripePaymentIntentId, timestamps. Record a \texttt{funding} transaction from
        the webhook, and a \texttt{release} transaction when the client marks the
        contract completed (this is the escrow ``release'' moment --- funds held on
        the platform account are now owed to the freelancer).
  \item Build \texttt{/freelancer/earnings}: available (released) balance, pending
        (funded, unreleased) balance, lifetime earned, transaction history table, and
        a monthly earnings chart. Data comes from real Transaction aggregation
        endpoints you add under \texttt{/api/earnings}.
  \item Build \texttt{/client/payments}: payment history per contract with receipts
        (amount, date, card brand/last4 from the PaymentIntent), plus total spent.
  \item Handle \texttt{charge.refunded} in the Stripe webhook: set contract
        \texttt{paymentStatus} to \texttt{refunded} and record a refund transaction.
  \item Add a clearly-labelled ``Withdraw'' UI on earnings that explains payouts are
        processed manually in this release (no Stripe Connect yet) and records a
        withdrawal request document for admin action. Update the Terms page so it
        describes exactly what the system now truly does.
  \item Pass the customer email to Stripe for receipts.
\end{enumerate}

\dod
\begin{itemize}[leftmargin=1.6em]
  \item Every payment produces a transaction row; earnings balances derive from
        transactions only; refund webhook verified with Stripe CLI.
\end{itemize}
\end{promptbox}

% ===========================================================================
\section{Phase 6 --- Communication \& Trust}
% ===========================================================================

\textit{Purpose: the human layer --- talking, being notified, and public
reputation.}

% ---------------------------------------------------------------------------
\objective{First-class messaging, the most-advertised missing feature.}

\begin{promptbox}{17}{Messaging System}
Implement real messaging between clients and freelancers.

\begin{enumerate}[leftmargin=1.6em]
  \item Models: \texttt{Conversation} (participants[2], optional contractId/
        projectId context, lastMessageAt, per-participant unread counts) and
        \texttt{Message} (conversationId, senderId, body, optional attachment URL,
        readAt). Indexed for inbox queries.
  \item APIs under \texttt{/api/messages}: list conversations, get/create a
        conversation with a specific user (idempotent), list messages (paginated,
        newest last), send message, mark conversation read. Participants only;
        validate everything with Zod; rate-limit sending.
  \item UI \texttt{/messages}: two-pane inbox (conversation list with unread badges
        and last-message preview; thread view with day separators, composer with
        Enter-to-send). Single-pane flow on mobile. Poll the open thread every few
        seconds and the inbox on an interval --- choose an approach that is simple
        and serverless-friendly; leave a clean seam to upgrade to SSE later.
  \item Entry points: ``Message'' buttons on freelancer profiles, gig pages, proposal
        cards, and the contract workspace all open the right conversation
        (creating it if absent). Wire the previously-disabled contact CTAs from
        Prompt 12.
  \item Unread message count badge on the app-shell sidebar/topbar; a
        \texttt{message} notification type for offline recipients (in-app +
        optional email respecting notification preferences).
\end{enumerate}

\dod
\begin{itemize}[leftmargin=1.6em]
  \item Two seeded users can hold a conversation across two browsers; unread badges
        update; every ``Message'' button in the product works.
\end{itemize}
\end{promptbox}

% ---------------------------------------------------------------------------
\objective{Complete the review loop in the UI and make reputation visible.}

\begin{promptbox}{18}{Reviews \& Reputation}
Build the review experience on top of the already-solid reviews API.

\begin{enumerate}[leftmargin=1.6em]
  \item After a contract completes, both parties see a prominent ``Leave a review''
        card on the contract workspace (and a notification link): star rating +
        comment, submitted to \texttt{POST /api/reviews}; show the mutual-review
        state (you reviewed / awaiting their review) from the contract detail API.
  \item Freelancer public profiles: full review section --- average, count, star
        breakdown bars, sorted list (recent/helpful toggle) with helpful-vote buttons
        wired to the existing endpoint. Client profiles get a lighter equivalent
        (rating shown on their project posts).
  \item Surface ratings consistently everywhere a person appears: freelancer cards,
        proposal cards, contract headers, project detail client card.
  \item Add a gentle reminder notification 3 days after completion if a party has
        not reviewed (a check performed on contract-page load or a simple scheduled
        route --- keep it serverless-friendly).
\end{enumerate}

\dod
\begin{itemize}[leftmargin=1.6em]
  \item Completing a contract and reviewing updates the target's public rating
        instantly; both directions work; helpful votes toggle.
\end{itemize}
\end{promptbox}

% ---------------------------------------------------------------------------
\objective{Notification coverage without gaps, plus email polish.}

\begin{promptbox}{19}{Notifications \& Email Completion}
Audit and complete the notification system end to end.

\begin{enumerate}[leftmargin=1.6em]
  \item Build the full-page inbox \texttt{/notifications} if Prompt 05's version was
        minimal: filters (all/unread), grouped by day, every item deep-linking to its
        subject, mark-all-read.
  \item Ensure EVERY significant event emits an in-app notification: proposal
        received, proposal accepted/rejected/withdrawn, contract funded, work
        delivered, contract completed/cancelled/disputed, payment succeeded/failed/
        refunded, new message, review received, admin role change. Grep every mutation
        route and fill the holes.
  \item Matching email coverage via \filepath{lib/email.ts} for the high-value subset
        (proposal accepted, contract funded, delivered, completed, payment events,
        review received), all respecting the user's notification preferences from
        Prompt 07. Convert inline awaited \texttt{sendEmail} calls in hot paths to
        fire-and-forget so slow email never delays API responses.
  \item Add missing transactional emails: verification email at credentials signup
        (token route + ``verify your email'' banner in-app until confirmed).
\end{enumerate}

\dod
\begin{itemize}[leftmargin=1.6em]
  \item A full hire-to-review loop generates the correct notification trail for both
        parties; credentials users can verify email.
\end{itemize}
\end{promptbox}

% ===========================================================================
\section{Phase 7 --- The Smart Layer (Gemini)}
% ===========================================================================

\textit{Purpose: the platform's namesake. Six AI pillars, built on one shared,
safe, server-side foundation. Never expose the API key to the browser; every
feature must degrade gracefully when AI is unavailable.}

% ---------------------------------------------------------------------------
\objective{One well-engineered AI foundation everything else plugs into.}

\begin{promptbox}{20}{AI Foundation --- Gemini Service Layer}
Create SkillSync's AI infrastructure using the Gemini API. No user-facing features
yet --- foundation only.

\begin{enumerate}[leftmargin=1.6em]
  \item Add \texttt{GEMINI\_API\_KEY} to the env schema and
        \filepath{.env.local.example}. Server-side only.
  \item Create \filepath{lib/ai/}: a singleton Gemini client; a
        \texttt{generateText} helper (system prompt + user content, JSON-mode option
        with Zod-validated output and one retry on parse failure); a
        \texttt{generateEmbedding} helper using Gemini's embedding model; a
        streaming helper for chat responses; central error handling that maps
        quota/network failures to a typed \texttt{AIUnavailableError}.
  \item Add an AI-specific rate limiter (per-user, e.g.\ 20 AI calls/hour) reusing
        \filepath{lib/rate-limit.ts} patterns, and response caching via
        \filepath{lib/cache.ts} keyed on input hashes for deterministic tasks.
  \item Create \texttt{POST /api/ai/health}: admin-only route that runs a trivial
        generation to confirm the integration.
  \item Add an \texttt{aiEnabled} helper: when the key is missing, every consumer
        must skip AI gracefully. Document the module in
        \filepath{docs/AI.md}: architecture, models used, cost controls, and the
        rule that AI never blocks a core flow.
\end{enumerate}

\dod
\begin{itemize}[leftmargin=1.6em]
  \item Health route returns a generation; removing the key degrades cleanly with no
        crashes; typecheck passes.
\end{itemize}
\end{promptbox}

% ---------------------------------------------------------------------------
\objective{Pillar 1: matching --- the platform's core intelligence.}

\begin{promptbox}{21}{Smart Matching Engine}
Build AI-powered matching between projects and freelancers.

\begin{enumerate}[leftmargin=1.6em]
  \item On project create/update, generate and store an embedding of
        title+description+skills on the Project document; on freelancer profile/gig
        update, store an embedding of headline+bio+skills+gig content on the User
        document. Backfill via a script for seeded data. (MongoDB Atlas Vector Search
        if available; otherwise store vectors and compute cosine similarity
        server-side over candidate sets prefiltered by category --- design the data
        access so the strategy can be swapped.)
  \item \texttt{GET /api/ai/matches/freelancers?projectId=} (project owner only):
        top freelancer matches --- similarity prefilter, then a Gemini rerank of the
        top $\sim$15 that returns for each a match percentage and TWO concrete
        reasons grounded in the actual profiles (JSON mode). Cache per
        project-version.
  \item \texttt{GET /api/ai/matches/projects} (freelancer): the inverse ---
        recommended open projects for the caller with reasons.
  \item UI: on the client's project detail page, a ``Smart Matches'' panel showing
        matched freelancer cards with the match percentage ring and reasons, plus
        Invite (opens messaging) and View Profile actions. On the freelancer
        dashboard, replace Prompt 13's placeholder ``Recommended for you'' with real
        AI matches carrying a ``Why this match?'' expander.
  \item If AI is unavailable, both panels fall back to category+skills overlap
        matching, visually identical minus the AI badge.
\end{enumerate}

\dod
\begin{itemize}[leftmargin=1.6em]
  \item A seeded React project surfaces the seeded React freelancers with plausible
        percentages and specific reasons; fallback path verified with the key
        removed.
\end{itemize}
\end{promptbox}

% ---------------------------------------------------------------------------
\objective{Pillars 3 \& 4: writing assistance and pricing intelligence.}

\begin{promptbox}{22}{AI Writing Studio \& Smart Pricing}
Embed generative assistance into the platform's three highest-friction writing
moments, plus pricing guidance.

\begin{enumerate}[leftmargin=1.6em]
  \item \textbf{Job-post composer} (client, in post-project): a ``Write with AI''
        panel where the client types a rough one-paragraph brief; Gemini returns a
        structured draft (title, description, suggested skills, suggested experience
        level) that PREFILLS the form for editing --- never auto-submits.
        Endpoint: \texttt{POST /api/ai/compose/project}.
  \item \textbf{Proposal assistant} (freelancer, in the proposal composer): given
        the project brief and the freelancer's profile, generate a tailored draft
        cover letter plus two talking-point suggestions; insert-into-editor UX.
        Endpoint: \texttt{POST /api/ai/compose/proposal}. Guard against prompt
        injection: treat project text strictly as data in your system prompt.
  \item \textbf{Profile optimiser} (freelancer, in settings/profile): reviews
        headline+bio and returns concrete rewrite suggestions and missing-skill
        prompts, shown side-by-side with one-click apply per suggestion.
        Endpoint: \texttt{POST /api/ai/compose/profile}.
  \item \textbf{Smart Pricing}: in post-project, after category+skills are chosen,
        show an AI budget estimate range with one-line rationale; in the proposal
        composer, show a suggested rate band derived from the project budget and the
        freelancer's rate. Endpoint: \texttt{POST /api/ai/pricing}.
  \item All endpoints: auth + role checks, the Phase-20 rate limiter, JSON-mode with
        Zod validation, and disabled-with-tooltip UI when AI is unavailable.
\end{enumerate}

\dod
\begin{itemize}[leftmargin=1.6em]
  \item All four assistants produce structured, editable output in their real forms;
        none can submit anything autonomously; rate limiting verified.
\end{itemize}
\end{promptbox}

% ---------------------------------------------------------------------------
\objective{Pillars 2 \& 5: decision support for hiring and trustworthy reviews.}

\begin{promptbox}{23}{Hiring Copilot \& Review Intelligence}
Build the client's AI decision support and review analysis.

\begin{enumerate}[leftmargin=1.6em]
  \item \textbf{Hiring Copilot}: on a project's proposals tab (3+ proposals), an
        ``Analyze proposals'' action calls \texttt{POST /api/ai/analyze/proposals}
        (owner-only) which returns, per proposal: a two-sentence summary, strengths,
        concerns (e.g.\ rate far outside budget, generic copy-paste letter), and an
        overall fit score --- rendered as a comparison view. Include an explicit
        on-screen disclaimer that this is assistance, not a decision; the ranking
        must never auto-reject anyone.
  \item \textbf{Review digest}: on freelancer profiles with 3+ reviews, a ``What
        clients say'' card --- Gemini summarises recurring themes (strengths +
        improvement areas) from review texts. Cached aggressively; regenerated only
        when new reviews arrive. Endpoint: \texttt{POST /api/ai/analyze/reviews}.
  \item \textbf{Anomaly flagging}: when a review is submitted, an async AI check
        scores it for spam/abuse/suspicious praise patterns; scores above threshold
        set a \texttt{flaggedForReview} field for the admin queue (Phase 8). Never
        block or delete automatically.
  \item Prompt-injection hardening for all three: user-generated text is data, never
        instructions; outputs are Zod-validated JSON.
\end{enumerate}

\dod
\begin{itemize}[leftmargin=1.6em]
  \item Copilot renders a useful comparison on seeded proposals; profile digest
        appears and caches; a gibberish review gets flagged in the DB.
\end{itemize}
\end{promptbox}

% ---------------------------------------------------------------------------
\objective{Pillar 6: the assistant that ties the smart experience together.}

\begin{promptbox}{24}{SyncMate --- The Platform Assistant}
Build SyncMate, a role-aware AI assistant available throughout the app shell.

\begin{enumerate}[leftmargin=1.6em]
  \item A floating assistant button in the app shell opening a chat panel; streaming
        responses via the Phase-20 streaming helper through
        \texttt{POST /api/ai/assistant} (auth required).
  \item System prompt gives SyncMate: the platform's actual feature set and routes,
        the user's role and first name, and strict boundaries (no legal/financial
        advice, no acting on the user's behalf, redirect abuse reports to support).
  \item Lightweight tool access, server-side: SyncMate can fetch the caller's own
        live context --- counts of their pending proposals, active contracts, unread
        messages --- to answer questions like ``what needs my attention?''. Read-only,
        strictly caller-scoped; the model must never receive other users' private
        data.
  \item Quick-start chips per role (client: ``How do I hire?'', ``Analyze my latest
        proposals''; freelancer: ``Find projects for me'', ``Improve my profile'') that
        deep-link to the relevant AI features where appropriate.
  \item Conversation history persists for the session (client-side is fine);
        rate-limited; hidden entirely when AI is unavailable.
\end{enumerate}

\dod
\begin{itemize}[leftmargin=1.6em]
  \item SyncMate answers platform questions with streaming text, reports the user's
        real pending counts, and refuses out-of-scope requests.
\end{itemize}
\end{promptbox}

% ===========================================================================
\section{Phase 8 --- Admin Panel}
% ===========================================================================

\objective{An operations cockpit: oversight, moderation, and dispute resolution.}

\begin{promptbox}{25}{Admin Panel}
Build the admin area at \texttt{app/(admin)/admin/\ldots}, middleware-gated to the
\texttt{admin} role, with its own dense two-pane shell.

\begin{enumerate}[leftmargin=1.6em]
  \item \textbf{Overview}: platform KPIs from \texttt{GET /api/admin/stats} extended
        with time series (signups, projects, GMV by month) rendered as charts.
  \item \textbf{Users}: the existing list/role/ban APIs get a real UI --- search,
        filters, detail drawer, role change and ban/unban with confirmations. THEN fix
        enforcement so banning matters: banned users are refused at login
        (credentials \texttt{authorize} and Google \texttt{signIn} callback), receive
        403s from authenticated APIs via a shared session guard, and are excluded
        from freelancer listings and public profiles --- not just search.
  \item \textbf{Content moderation}: admin endpoints + UI to view and unpublish/
        remove projects, gigs, and reviews; a queue for reviews auto-flagged by
        Prompt 23, with approve/remove actions.
  \item \textbf{Disputes}: queue of disputed contracts showing both parties, the
        dispute reason, payment state, and linked messages; admin resolution actions
        --- resolve to completed (release) or cancelled (mark for refund) --- with an
        audit note stored on the contract and notifications to both parties.
  \item \textbf{Operations}: contact-form inbox (list \texttt{Contact} documents,
        mark \texttt{readAt}); withdrawal-request queue from Prompt 16 with
        mark-as-paid.
  \item Fix the stale-JWT problem: on role change or ban, force re-validation ---
        re-check the user document inside the NextAuth JWT callback on an interval
        (e.g.\ refresh role/ban claims from the DB when the token is older than
        15 minutes).
\end{enumerate}

\dod
\begin{itemize}[leftmargin=1.6em]
  \item Non-admins cannot reach \texttt{/admin} or its APIs; banning a user locks
        them out within one token refresh; a dispute can be resolved end to end.
\end{itemize}
\end{promptbox}

% ===========================================================================
\section{Phase 9 --- Hardening \& Launch Quality}
% ===========================================================================

\objective{Close the security gaps found in the audit, systematically.}

\begin{promptbox}{26}{Security Hardening Pass}
Execute a focused security pass over the entire API surface.

\begin{enumerate}[leftmargin=1.6em]
  \item \textbf{Rate limiting}: add sensible limits to forgot-password,
        reset-password, credentials login attempts (with temporary lockout after
        repeated failures per email+IP), proposals, reviews, messages, uploads, and
        all AI routes. Centralise limits in one config.
  \item \textbf{Input integrity}: a shared \texttt{isValidObjectId} guard returning
        400 on every \texttt{[id]} route (users, projects, skills, proposals,
        payments currently 500); escape user input in every remaining
        \texttt{\$regex} construction; cap the number of files per upload request and
        add per-user daily upload quotas.
  \item \textbf{AuthZ sweep}: re-verify every mutation route enforces role and
        ownership after all phases; write the result as a route-by-route
        authorization table in \filepath{docs/SECURITY.md}.
  \item \textbf{Secrets \& env}: add build-time env validation (typed env module)
        covering every variable including \texttt{GEMINI\_API\_KEY}; confirm nothing
        secret is exposed with \texttt{NEXT\_PUBLIC\_}.
  \item \textbf{Webhook \& payments}: verify webhook idempotency (duplicate event
        delivery must be safe); confirm refund/dispute events are handled.
  \item Clean up remaining validation-schema drift by making \filepath{types/schemas.ts}
        the single source of truth imported by routes.
\end{enumerate}

\dod
\begin{itemize}[leftmargin=1.6em]
  \item Malformed IDs return 400 everywhere; brute-force login is throttled;
        \filepath{docs/SECURITY.md} authz table is complete and accurate.
\end{itemize}
\end{promptbox}

% ---------------------------------------------------------------------------
\objective{The final polish and a regression safety net.}

\begin{promptbox}{27}{Quality Pass, SEO \& Test Foundation}
Final production-readiness pass.

\begin{enumerate}[leftmargin=1.6em]
  \item \textbf{UX audit}: every async view has loading skeletons, empty states
        with a helpful CTA, and error states with retry. Audit at 375px, 768px and
        1440px; fix responsive breaks. Verify dark mode on every page.
  \item \textbf{SEO}: per-page \texttt{metadata} (title/description) across
        marketing, gig and freelancer-profile pages; Open Graph tags; sitemap and
        robots for public routes only; \texttt{noindex} on the app shell.
  \item \textbf{Accessibility}: keyboard navigation through the main flows,
        focus-visible states, aria labels on icon buttons, form error associations,
        colour-contrast check on both themes.
  \item \textbf{Performance}: eliminate client components that can be server
        components; verify no oversized bundles (dynamic-import heavy widgets like
        charts and the AI panels).
  \item \textbf{Tests}: install Vitest; unit-test the critical money/state logic
        (contract transition rules, proposal accept transaction, transaction ledger
        math, rate-limit behaviour, AI JSON-parsing fallbacks). Install Playwright;
        one smoke E2E covering register $\rightarrow$ onboard $\rightarrow$ post
        project $\rightarrow$ bid $\rightarrow$ accept $\rightarrow$ fund (Stripe
        test mode) $\rightarrow$ deliver $\rightarrow$ complete $\rightarrow$ review.
        Add \texttt{npm test} and a CI-ready script.
  \item Update the README and \filepath{docs/} to describe the platform as it now
        actually is, including an env-vars table and the AI feature overview.
\end{enumerate}

\dod
\begin{itemize}[leftmargin=1.6em]
  \item The E2E smoke passes clean; unit tests cover the money paths; docs match
        reality.
\end{itemize}
\end{promptbox}

% ===========================================================================
\section{Appendix A --- Environment Variables After All Phases}
% ===========================================================================

\begin{table}[h]
\small
\begin{tabularx}{\textwidth}{@{}lX@{}}
\toprule
\textbf{Variable} & \textbf{Purpose} \\
\midrule
\texttt{MONGODB\_URI}                        & MongoDB Atlas connection string \\
\texttt{NEXTAUTH\_SECRET} / \texttt{NEXTAUTH\_URL} & Auth.js signing secret and canonical URL \\
\texttt{GOOGLE\_CLIENT\_ID} / \texttt{\_SECRET} & Google OAuth \\
\texttt{STRIPE\_SECRET\_KEY}                 & Stripe server key \\
\texttt{STRIPE\_WEBHOOK\_SECRET}             & Webhook signature verification \\
\texttt{NEXT\_PUBLIC\_STRIPE\_PUBLISHABLE\_KEY} & Stripe.js in the browser \\
\texttt{CLOUDINARY\_*}                       & File and avatar storage \\
\texttt{RESEND\_API\_KEY} / \texttt{RESEND\_FROM} & Transactional email \\
\texttt{GEMINI\_API\_KEY}                    & All Smart-layer features (server-only) \\
\texttt{NEXT\_PUBLIC\_APP\_URL}              & Absolute URL building \\
\bottomrule
\end{tabularx}
\end{table}

% ===========================================================================
\section{Appendix B --- Vibe-Coding Discipline}
% ===========================================================================

\begin{itemize}[leftmargin=1.6em]
  \item \textbf{Paste the prompt verbatim first}, then answer the agent's clarifying
        questions from this document rather than improvising new scope.
  \item \textbf{Demand proof.} After each prompt ask the agent to demonstrate the
        Definition-of-done items concretely (show the DB document, show the network
        call, show the redirect) --- this codebase previously ``passed'' on
        \texttt{setTimeout} theatre.
  \item \textbf{Commit per prompt} with the prompt number in the message; branch per
        phase if you want extra safety.
  \item \textbf{When a prompt is too big for one session}, split it at the numbered
        item boundaries --- the internal ordering within each prompt is also
        dependency-ordered.
  \item \textbf{Keep the seed script current.} Whenever a model changes, the same
        session must update \filepath{scripts/seed.ts}; a stale seed silently breaks
        every later verification.
\end{itemize}

\vspace{1cm}
\begin{center}
  \color{phasegray}\itshape
  End of the Prompt Diary. Execute in order --- and make ``Smart'' mean something.
\end{center}

\end{document}
